\chapter{Formation, Clustering and Coordination}
\label{ch:formation}

The distributed estimator and controller of the previous two chapters provide
each agent with a state estimate and a safe input. This chapter defines the
cluster-level behaviours those primitives serve---how a formation is specified,
how the graph is maintained, and how the cluster reconfigures and tolerates
faults.

\section{Formation specifications}

Three standard specifications of a formation are supported, in increasing order
of decentralisation \cite{oh2015survey}.
\begin{itemize}[leftmargin=1.4em]
  \item \textbf{Position-based (virtual structure).} The formation is a rigid
        body---the ``virtual structure''---and each agent tracks a fixed offset
        $\vv{d}_i$ from its moving frame \cite{ren2004formation}. This is the
        specification used in \cref{ch:dmpc}: the centre $\vv{c}(t)$ is the
        virtual structure's origin, agreed by consensus, and $\vv{d}_i$ places
        agent $i$. It is simple and exact but requires a common reference frame.
  \item \textbf{Displacement-based (leader--follower).} Followers maintain
        prescribed relative displacements to a leader (or to designated
        neighbours) in a common orientation. The fixed-wing study of
        \cref{ch:simulation} uses this form: a virtual leader traverses a loiter
        and followers hold along-track and cross-track offsets.
  \item \textbf{Distance-based (rigidity).} Only inter-agent distances are
        prescribed; the formation is maintained if the underlying framework is
        infinitesimally rigid \cite{krick2009stabilisation}. This is the most
        decentralised---no common frame is needed---and the most subtle, since
        distance constraints admit reflections and flips.
\end{itemize}
In all three the desired geometry enters the controller as the reference slot or
as inter-agent constraints, and the DMHE supplies the relative-position estimates
that displacement- and distance-based schemes require.

\section{The communication graph and its maintenance}

The graph $\Graph$ of \cref{sec:graph} is not static: links appear and vanish as
agents move in and out of communication range, and the topology must stay
connected for consensus and coordination to function. Two mechanisms keep it so.

\paragraph{Neighbour discovery.} Each agent periodically broadcasts a beacon with
its identity, estimated position and covariance; agents within range that hear it
add the sender to $\Neigh_i$ and drop neighbours whose beacons lapse. The
resulting time-varying graph is exactly the switching topology whose consensus
convergence is guaranteed, under a joint-connectivity condition, by
Olfati-Saber and Murray \cite{olfati2004consensus,olfati2007consensus} and by Fax
and Murray \cite{faxmurray2004}.

\paragraph{Connectivity maintenance.} Because the algebraic connectivity
$\lambda_2(\Lap)$ sets the coordination rate and must stay positive, the cluster
monitors it (each agent can estimate $\lambda_2$ by a distributed power
iteration on $\Lap$) and the controller penalises actions that would break a
critical link---equivalently, a connectivity-preserving CBF of the form
\cref{eq:cbf-h} with $h=R_{\text{comm}}^2-\norm{\vv{p}_i-\vv{p}_j}^2$ keeps a
needed neighbour within range. Collision avoidance (a lower bound on separation)
and connectivity maintenance (an upper bound on the distance to a required
neighbour) are thus the same mechanism with opposite sign.

\section{Reconfiguration and fault handling}

\paragraph{Reconfiguration.} A formation change is a change of the offsets
$\{\vv{d}_i\}$ (or of the prescribed distances). Because the offsets enter only
through the MPC reference and the coupling constraints, a reconfiguration is
executed simply by ramping $\vv{d}_i$ to their new values; the controllers track
the moving slots and the safety filter of \cref{sec:cbf} guarantees the transient
stays collision-free even when the intermediate geometry brings agents close.
The simulation study exercises exactly this: a six-agent ring contracts and
rotates mid-mission.

\paragraph{Fault handling.} Agent faults are detected through the same beacon
mechanism: a lapsed beacon, or an innovation/residual monitor on a neighbour's
broadcast estimate, flags a failed or misbehaving agent. On detection the cluster
\emph{re-clusters}: the failed agent is removed from every $\Neigh_i$, the
formation offsets are re-assigned among the survivors (a small assignment problem
solved by consensus), and the virtual structure is re-scaled. The over-actuation
of the hexacopter (\cref{sec:hexa-mixer}) provides a second, vehicle-level layer
of fault tolerance: an agent that loses a single rotor remains controllable in
thrust and tilt and can hold or leave the formation gracefully rather than
falling from it \cite{michieletto2018fundamental}.

\paragraph{Heterogeneity.} Because the coordination layer sees each agent only
through the common interface of \cref{sec:capi}---a state estimate in, a safe
input out---a cluster may mix hexacopters and fixed-wing aircraft freely. The
vehicle class changes $\vv{f}_i$, the constraint set $\vv{c}_i$ (a fixed-wing
carries a minimum-airspeed row a hexacopter does not) and the achievable
formation geometry (fixed-wing slots must respect turn-rate limits), but not the
distributed algorithms of \cref{ch:dmhe,ch:dmpc}. This is the payoff of routing
all vehicle-specific detail through the per-agent solvers.
